\section{Introduction}\label{sec:introduction}

Ramsey's theorem guarantees that every graph on $N\ge4^{K-1}$ vertices
contains a clique or an independent set of order $K$.  Its classical
constructive proof \cite{ErdosSzekeres1935} repeatedly chooses a pivot and
restricts the candidate set to one of the pivot's two colour
neighbourhoods; truncating each retained neighbourhood to half its current
size, it finds a homogeneous $K$-set with $O(N)$ adjacency queries.  The
corresponding total search problem, \textsc{Ramsey}, has been studied in
proof complexity and in TFNP
\cite{Krajicek2001,Krajicek2005,KomargodskiNaorYogev2019,
PasarkarPapadimitriouYannakakis2023,JainLiRobereXun2024}: the input graph on
$N=2^n$ vertices is given by a circuit or an oracle, and the target order is
$K=n/2$.  The query complexity of this relation in the deterministic,
randomized, and quantum models is posed as an open question
in \cite[Section~III]{JainLiRobereXun2024}.  This paper answers the
quantum question up to a polynomial gap and clarifies the lower-bound
record on both the quantum and the classical side.

The central observation is that the candidate sets of the constructive proof
need not be materialized.  After $i$ rounds, membership in the current set
is determined by the $i$ adjacency constraints imposed by the previous
pivots.  These short predicates allow quantum search to sample from the
same recursion at substantially smaller query cost.  The threshold $4^{K-1}$
is the natural scale for this recursive construction.

\subsection{Our results}

We work in the promised valid-graph adjacency-oracle model.  One coherent
query returns the colour of a requested unordered vertex pair, and the output
is a set of vertices inducing one colour.  Our main result is the following.

\begin{table}[t]
\centering
\setlength{\tabcolsep}{4pt}
\begin{tabular}{llll}
\hline
model & input & lower bound & upper bound\\
\hline
randomized & worst case
 & $\Omega(N^{1-1/\sqrt2})$ \ [\cref{prop:classical-lower}]
 & $O(N)$ \ \cite{ErdosSzekeres1935}\\
randomized & $G(N,1/2)$
 & $\Omega(N^{1-1/\sqrt2})$ \ [\cref{prop:classical-lower}]
 & $\widetilde O(N^{1/2})$ \ [\cref{rem:classical-greedy}]\\
quantum & worst case
 & $\Omega(N^{1/12})$ \ [\cref{prop:quantum-lower}]
 & $\widetilde O(N^{1/2})$ \ [\cref{cor:succinct}]\\
quantum & $G(N,1/2)$
 & open
 & $\widetilde O(N^{1/4})$ \ [\cref{cor:separation}]\\
\hline
\end{tabular}
\caption{Bounded-error query complexity of finding a homogeneous set of
order $\lfloor n/2\rfloor+1$ in a graph on $N=2^n$ vertices, in the worst
case and on the uniform distribution $G(N,1/2)$; here
$1-1/\sqrt2\approx0.2929$.  The lower bounds are the random-Painter
estimate of \cite{ConlonFoxGrinshpunHe2019} and the collision reduction of
\cite{JainLiRobereXun2024} at multiplicity two (\cref{prop:quantum-lower}
is stated for order $\lfloor n/2\rfloor$ and hence valid here).}
\label{tab:bounds}
\end{table}

\begin{theorem}[Implicit-majority quantum Ramsey search]
\label{thm:main}
Let $K\ge2$, $N\ge4^{K-1}$, and $0<\eta<1/2$.  Given coherent edge-oracle
access to a simple undirected graph on $N$ vertices, there is a quantum
algorithm which, with probability at least $1-\eta$, outputs a $K$-clique or
a $K$-vertex independent set using
\[
 O\!\left(2^K K\log\frac K\eta\right)
\]
edge queries in the worst case.  The algorithm verifies every returned
witness and may output $\bot$ on the exceptional event.
\end{theorem}

At the succinct Ramsey scale, the theorem gives the following form.

\begin{corollary}\label{cor:succinct}
Let $n\ge2$, $N=2^n$, and $0<\eta<1/2$.  Given coherent edge-oracle access
to a simple graph on $N$ vertices, one can find a clique or independent set
of order $\lfloor n/2\rfloor+1$ using
\[
 O\!\left(
 \sqrt N\log N\log\frac{\log N}{\eta}
 \right)
\]
edge queries in the worst case.
\end{corollary}

\begin{proof}
Set $K=\lfloor n/2\rfloor+1$.  Then
$4^{K-1}=2^{2\lfloor n/2\rfloor}\le 2^n=N$, so \cref{thm:main} applies, and
$2^K\le2\sqrt N$, $K=O(\log N)$, and
$\log(K/\eta)=O(\log(\log N/\eta))$.
\end{proof}

To our knowledge, this is the first sublinear worst-case query algorithm
for the Ramsey relation in any model.  \Cref{tab:bounds} places it among
the known bounds together with the further results of this paper.

\paragraph{An estimation-free recursion.}
Following the colour of a uniformly sampled edge selects a colour class with
probability proportional to its size, and a one-dimensional survival
recurrence shows that this recursion completes with probability greater
than $1/2$ from $4^{K-1}$ vertices.  This gives an
$O(2^KK^2\log K\log(1/\eta))$-query algorithm without any majority
estimation (\cref{prop:size-biased-algorithm}), which extends to every
fixed number of edge colours (\cref{cor:multicolour}).

\paragraph{A randomized lower bound.}
Every randomized algorithm needs $\Omega(M^{1-1/\sqrt2})$ queries to find a
homogeneous $K$-set in a graph on $M=4^{K-1}$ vertices, even on the uniform
distribution $G(M,1/2)$ (\cref{prop:classical-lower}); this improves the
$M^{1/4}$ black-box bound of
\cite{ImpagliazzoNaor1988,KomargodskiNaorYogev2019}.  The proof transports
the random-Painter weight estimate of \cite{ConlonFoxGrinshpunHe2019} to
the query model, with a self-contained version in \cref{app:classical-lower}.

\paragraph{A quantum speedup on random graphs.}
On $G(N,1/2)$, a greedy quantum neighbourhood search finds a $K$-clique
with $\widetilde O(N^{1/4})$ queries (\cref{prop:random-graph}).  Together
with the distributional form of the lower bound, this gives a polynomial
separation between the quantum and randomized query complexities of the
Ramsey relation on its natural hard distribution (\cref{cor:separation}).
We are not aware of an earlier provable quantum speedup for this relation.
In the worst case, no separation is claimed: the randomized complexity is
only known to lie between $N^{1-1/\sqrt2}$ and $N$.

\paragraph{The quantum lower-bound record.}
The paragraph of \cite{JainLiRobereXun2024} that raises the query question
also reports an $N^{1-o(1)}$ quantum lower bound at the default parameters,
obtained by composing its collision reduction with the multicollision lower
bound of \cite{LiuZhandry2019}.  \Cref{thm:main} shows that no such bound
can hold at these parameters, and \cref{sec:jlrx} shows that the cited
ingredients compose to an exponent of at most $1/12$ (at most $1/24$ under
the parameter condition as printed), attained at multiplicity two.  The
resulting $\Omega(N^{1/12})$ bound (\cref{prop:quantum-lower}) is, to our
knowledge, the best known quantum lower bound; the main theorems of
\cite{JainLiRobereXun2024} do not depend on the remark.

\subsection{Techniques}

Both algorithms represent the candidate set after round $i$ by the
conjunction of the adjacency constraints generated by the first $i$ pivots,
so a membership test costs $O(i)$ edge queries, and sample it with capped
unknown-solution search (\cref{lem:capped-search}); conditional on success,
permutation symmetry makes each sample exactly uniform.  The scale-aware
recursion estimates a near-majority colour at every pivot.  Sampling a deep
set is exponentially more expensive than sampling an early one, so we
distribute a fixed error budget across the levels, doubling the allowed
error every six levels; a product invariant then keeps every candidate set
within a constant factor of its ideal size, and the query sum is dominated
by the last levels.

The lower bound and the random-graph speedup both rest on the predictable
size of common neighbourhoods in $G(N,1/2)$.  A martingale weight argument
in the style of \cite{ConlonFoxGrinshpunHe2019} shows that an adaptive
algorithm with $T$ queries exposes a monochromatic $K$-set with probability
at most $2\cdot2^{-\binom K2+c(c-1)}(2T)^{K-c}$ for every $c\le K/2$, and
optimizing $c$ gives the exponent $2-\sqrt2$; a union bound shows that
every $j$-set has common neighbourhood of density about $2^{-j}$, so greedy
quantum search costs $O(2^{j/2})$ predicate evaluations at depth $j$.

\paragraph{Organization.}
\Cref{sec:model} fixes the query model and the sampler.
\Cref{sec:algorithm} presents the two recursions, their analysis, and the
proof sketch of \cref{thm:main}; \cref{sec:lower} states the lower bounds
and the random-graph speedup.  \Cref{sec:prior-art} discusses related work
and the reported lower bound, and \cref{sec:scope} lists open questions.
All full proofs are in the appendices: the sampler and the encoding
corollary in \cref{app:capped-search}, the upper bounds in
\cref{app:upper-proofs}, the lower bounds and random-graph results in
\cref{app:classical-lower}, and the multicolour extension, numerical
illustrations, and AI disclosure in \cref{app:extensions}.
